\documentclass[a4paper]{article}
\usepackage[pdftex]{graphicx}
\usepackage[utf8]{inputenc}
\usepackage{enumerate}
\usepackage{amssymb}
\usepackage{icomma}
\usepackage{href-ul}
\hypersetup{
	colorlinks=true,
	linkcolor=blue,
	urlcolor=blue}
\usepackage{siunitx}
\sisetup{locale=DE} 
\usepackage{geometry}
\geometry{a4paper, top=15mm, left=15mm, right=15mm, bottom=15mm,
	headsep=10mm, footskip=12mm}

\begin{document}
{\bf Dichte, Volumen und Masse}

\begin{enumerate}[1.]
{\bf \item Fülle die Tabelle aus.} Beachte: die Dichte ist für jeden Stoff konstant. Das Volumen und die entsprechende Masse sind beliebig wählbar. Der Zustand ist entweder fest oder flüssig oder gasförmig. 

\renewcommand{\arraystretch}{2}
\begin{tabular}{|p{1.7 cm}|p{1.7 cm}|p{1.7 cm}|p{1.7 cm}|p{1.7 cm}|p{1.7 cm}|p{1.7 cm}|p{1.7 cm}|}
\hline
Stoff	& Holz & Blei & Wasser & Luft & Speiseöl & Gold & Aluminium \\
\hline 
Zustand &  &  &  &  &  &  &  \\
\hline
Beispiel & Bauklotz & Anglerblei & Flasche & Ballon & Flasche & Barren & Spitzer\\
\hline
Masse & $\SI{35}{\gram}$  & $\SI{5,0}{\gram}$  &  & $\SI{5,2}{\gram}$  & $\SI{0,63}{\kilogram}$  &  &$\SI{4,3}{\gram}$  \\
\hline
Volumen &$\SI{50}{\centi\meter^3}$  &  &$\SI{1,5}{\deci\meter^3}$  &  & $\SI{0.7}{\deci\meter^3 }$  &$\SI{0,64}{\deci\meter^3}$  & $\SI{1,6}{\centi\meter^3}$  \\
\hline
Dichte & &$\SI{11,3}{\gram\over{\centi\meter^3}}$  &$\SI{1,0}{\gram\over{\centi\meter^3}}$  &$\SI{1.3}{\gram\over{\deci\meter^3}}$  &   &$\SI{19,3}{\gram\over{\centi\meter^3}}$  &  \\
\hline
\end{tabular}
\vspace{0.5 cm}

\begin{tabular}{|p{1.4 cm}|p{1.4 cm}|p{1.5 cm}|p{1.8 cm}|p{1.7 cm}|p{2.0 cm}|p{2.1 cm}|p{1.7 cm}|}
\hline
Stoff	& Kork & Eisen & Styropor & Wasserstoff & Iridium & Quecksilber & Polyethylen \\
\hline
Zustand &  &  &  &  & fest &  &  \\
\hline 
Beispiel & Korken & Schlüssel & Dämmplatte & Ballon & Urkilogramm & Thermometer & Spielwürfel\\
\hline
Masse &$\SI{1,9}{\gram}$  &  & $\SI{150}{\gram}$  & $\SI{0,41}{\gram}$  &  & $\SI{0,67}{\gram}$ & $\SI{3,1}{\gram}$\\
\hline
Volumen & & $\SI{1,1}{\centi\meter^3}$  & $\SI{5,0}{\deci\meter^3}$  &  & $\SI{45}{\centi\meter^3}$  &$\SI{0,05}{\centi\meter^3}$  & \\
\hline
Dichte &$\SI{0,3}{\gram\over{\centi\meter^3}}$  &$\SI{7,7}{\gram\over{\centi\meter^3}}$  &  &$\SI{90}{\gram\over{\meter^3}}$  &$\SI{22,4}{\gram\over{\centi\meter^3}}$  &  &$\SI{0,95}{\gram\over{\centi\meter^3}}$  \\
\hline
\end{tabular}

{\bf \item Welches ist der spezifisch leichteste Stoff ?} Welches ist der mit der höchsten Dichte ?

{\bf \item Warum fliegt ein Wasserstoffballon nach oben?} Sein Gewicht ist doch nicht negativ...

{\bf \item Aufgabe}

\begin{enumerate}[a)]
\item Welcher der obengenannten festen Stoffe  würde in welcher Flüssigkeit schwimmen können?

\renewcommand{\arraystretch}{4}
\begin{tabular}{|p{3 cm}|p{12.5 cm}|}
	\hline
Schwimmt in & Stoff \\
\hline
Speiseöl & \\
\hline
Wasser	& \\
\hline
Quecksilber & \\
\hline
\end{tabular}

\item Welcher Stoff schwimmt in Wasser, geht aber in Speiseöl unter?
\item Welche Stoffe gehen in Quecksilber unter ?
\end{enumerate} 
\end{enumerate} 
\fbox{
	\begin{minipage}{0.5\textwidth}
		Zur Lösung bitte \href{https://www.okuyakl.de/phys/p7dihL073/ll073.pdf}{hier klicken} oder den QR-Code scannen.\\
	Weitere Arbeitsblätter gibt es unter 
	
	\href{https://www.okuyakl.de}{www.okuyakl.de}
	\end{minipage}
	\hfill
	\begin{minipage}{0.4\textwidth}
		\includegraphics[width=1.5 cm]{../../viecher/zwe03}
		\includegraphics[width=3 cm]{qr073}
		\includegraphics[width=2 cm]{../../viecher/afanticon1}
		
\end{minipage}}

\end{document}

